\documentclass[11pt,a4paper]{article}

% ── Packages ──────────────────────────────────────────────────────────────────
\usepackage[margin=2.5cm]{geometry}
\usepackage{amsmath,amssymb}
\usepackage{booktabs}
\usepackage{array}
\usepackage{listings}
\usepackage{xcolor}
\usepackage{hyperref}
\usepackage{parskip}
\usepackage{titlesec}
\usepackage{enumitem}
\usepackage{microtype}

% ── Colours ───────────────────────────────────────────────────────────────────
\definecolor{codebg}{RGB}{245,245,245}
\definecolor{codeframe}{RGB}{200,200,200}
\definecolor{keywordcolor}{RGB}{0,100,160}
\definecolor{commentcolor}{RGB}{100,120,100}
\definecolor{stringcolor}{RGB}{160,60,0}

% ── Code listing style ────────────────────────────────────────────────────────
\lstdefinestyle{bashstyle}{
  language=bash,
  backgroundcolor=\color{codebg},
  frame=single,
  framerule=0.4pt,
  rulecolor=\color{codeframe},
  basicstyle=\ttfamily\small,
  keywordstyle=\color{keywordcolor}\bfseries,
  commentstyle=\color{commentcolor}\itshape,
  stringstyle=\color{stringcolor},
  breaklines=true,
  breakatwhitespace=false,
  showstringspaces=false,
  tabsize=2,
  captionpos=b,
}

\lstdefinestyle{pythonstyle}{
  language=Python,
  backgroundcolor=\color{codebg},
  frame=single,
  framerule=0.4pt,
  rulecolor=\color{codeframe},
  basicstyle=\ttfamily\small,
  keywordstyle=\color{keywordcolor}\bfseries,
  commentstyle=\color{commentcolor}\itshape,
  stringstyle=\color{stringcolor},
  breaklines=true,
  showstringspaces=false,
  tabsize=4,
  captionpos=b,
}

\lstdefinestyle{yamlstyle}{
  backgroundcolor=\color{codebg},
  frame=single,
  framerule=0.4pt,
  rulecolor=\color{codeframe},
  basicstyle=\ttfamily\small,
  breaklines=true,
  showstringspaces=false,
  tabsize=2,
}

% ── Hyperref setup ────────────────────────────────────────────────────────────
\hypersetup{
  colorlinks=true,
  linkcolor=blue!60!black,
  citecolor=green!50!black,
  urlcolor=blue!70!black,
}

% ── Section formatting ────────────────────────────────────────────────────────
\titleformat{\section}{\large\bfseries}{}{0em}{}[\titlerule]
\titleformat{\subsection}{\normalsize\bfseries}{}{0em}{}

% ── Inline code helper ────────────────────────────────────────────────────────
\newcommand{\ic}[1]{\texttt{#1}}

% ──────────────────────────────────────────────────────────────────────────────
\begin{document}

\title{CHEETAH Systems Notes}
\author{}
\date{Tuesday, May 12, 2026}
\maketitle
% \tableofcontents
% \newpage
% ──────────────────────────────────────────────────────────────────────────────
\subsection{Pipeline to Run Experiments}

\begin{enumerate}
  \item \textbf{\ic{run\_pipeline\_timed.sh}} is the entry point. This fully
        runs 3 trajectories, N iterations each:
        \begin{itemize}[noitemsep]
          \item Vizier + Timeloop + FAST original mini-LP
          \item Vizier + Timeloop + V0-LP
          \item Vizier + V1-LP
        \end{itemize}

  \item \textbf{\ic{timeloop\_to\_corgi.py}} deals with the CORGIsolver details. Selecting the wrong config/default solver crashes the machine. CORGIsolver mitigates this problem.
\end{enumerate}

% ══════════════════════════════════════════════════════════════════════════════
\subsection{Process-Normalized TPU-v3 Baseline and Perf/TDP}
% ══════════════════════════════════════════════════════════════════════════════

Three methods: Static, V0, V1.

% ──────────────────────────────────────────────────────────────────────────────
\subsection{1. Sub-10-nm Normalised TPU-v3 Baseline}

From the FAST paper (Section~6.1, ``Methodology and Simulator''):

\begin{quote}
\textit{``To estimate area and power consumption, we built analytical models
correlated to production designs on an industry sub-10-nm manufacturing
process; for a fair comparison, the TPU-v3 baseline is also modelled using this
same process technology. TDP is estimated as power virus power.''}
\end{quote}

Figure~10's caption is  more clear:

\begin{quote}
\textit{``Perf/TDP numbers relative to a hypothetical TPU-v3 die-shrunk to the
same sub-10-nm process technology.''}
\end{quote}

So in CHEETAH we assume ``TPU-v3 baseline normalized to sub-10-nm'' as an one-to-one adjustment to prevent candidate designs from winning because of a node-shrink instead of scheduling. Since Vizier = outer Bayesian loop, the cost models live inside the scoring function. 
% ──────────────────────────────────────────────────────────────────────────────
\subsection{2. Where $8.2 \pm 2.7\,\%$ Comes From}

We inherit this number. FAST measured it by comparing the
simulator-reported QPS against profiled QPS from real TPU-v3. We use
the same Vizier + Timeloop + TPU-v3 baseline configuration, and match calibration to inherit this (additional checks against TPUv3 white paper cited).

% ──────────────────────────────────────────────────────────────────────────────
\subsection{3. Perf/TDP Calculations}

We use the equation:

\begin{equation}
\text{Perf/TDP}(\times)
= \frac{\text{QPS}_{\text{design}} / \text{TDP}_{\text{design}}}
       {\text{QPS}_{\text{TPU-v3}} / \text{TDP}_{\text{TPU-v3}}}
= \frac{\text{Speedup}_{\text{vs TPU-v3}}}{\text{TDP ratio}}
\end{equation}

\subsubsection*{Speedup Numerator}

\begin{equation}
\text{Speedup}_{\text{vs TPU-v3}}
= \frac{\text{QPS}_{\text{design}}}{\text{QPS}_{\text{TPU-v3}}}
= \frac{\text{cycles}_{\text{TPU-v3}}}{\text{cycles}_{\text{design}}}
\end{equation}

\begin{itemize}[noitemsep]
  \item \ic{cycles\_design} = the LP objective at the candidate HW config.
  \item \ic{cycles\_TPU-v3} = the LP objective at TPU-v3 HW.
\end{itemize}

\subsubsection*{TDP Ratio Denominator}

The TDP analytical model collapses to a linear two-term proxy with a 0.30
floor:

\begin{equation}
\text{TDP ratio}
= \max\!\left(
0.70 \cdot \frac{\text{MACs}_{\text{design}}}{\text{MACs}_{\text{TPU-v3}}}
+ 0.30 \cdot \frac{\text{BW}_{\text{design}}}{\text{BW}_{\text{TPU-v3}}},
\ 0.30
\right)
\end{equation}

\begin{table}[h]
\centering
\begin{tabular}{lll}
\toprule
\textbf{Symbol} & \textbf{Value} & \textbf{Source} \\
\midrule
\ic{MACs\_TPU-v3} & 65,536 & $256 \times 256$ systolic array \\
\ic{BW\_TPU-v3}   & 900\,GB/s & spec \\
$\alpha = 0.70$ & compute share of TDP & calibrated to FAST model \\
$\beta  = 0.30$ & memory share of TDP  & same \\
$\tau   = 0.30$ & TDP floor            & prevents tiny-chip \\
\bottomrule
\end{tabular}
\caption{TDP proxy constants.}
\end{table}

\ic{MACs\_design} is computed from the Vizier-proposed HW dimensions:
\begin{equation}
\text{MACs}_{\text{design}}
= \text{pe\_rows} \times \text{pe\_cols} \times \text{systolic\_x} \times \text{systolic\_y}
\end{equation}

\ic{BW\_design} is the \ic{dram\_bandwidth\_gbps} Vizier dimension.

\subsubsection*{The Function Implementation}

\begin{lstlisting}[style=pythonstyle]
TPU_V3_MACS    = 65_536     # 256 x 256 systolic
TPU_V3_BW_GBPS = 900        # HBM2

def perf_tdp(cycles_design, cycles_tpu_v3, macs_design, bw_design):
    speedup   = cycles_tpu_v3 / cycles_design
    tdp_ratio = max(
        0.70 * (macs_design / TPU_V3_MACS)
        + 0.30 * (bw_design  / TPU_V3_BW_GBPS),
        0.30,
    )
    return speedup / tdp_ratio
\end{lstlisting}

% ──────────────────────────────────────────────────────────────────────────────
% \subsection{4. Per Pipeline}

The Perf/TDP equation is identical for both pipelines, how cycles\_design is produced differs since V1 takes Timeloop out of the loop. 

% % ──────────────────────────────────────────────────────────────────────────────
% \subsection{5. Worked Example --- A100 + V1 LP on Llama-3 70B}

% \begin{lstlisting}[style=pythonstyle]
% cycles_design   =  29_527_900_160     # V1 LP solve at A100 HW
% cycles_TPU_v3   = 226_410_124_872     # TPU-v3 baseline for llama3_70b
% speedup         = 226_410_124_872 / 29_527_900_160   # = 7.67x

% # A100 hardware knobs:
% #   MACs_design = 163_840   BW_design = 2039 GB/s

% tdp_ratio = max( 0.70 * (163_840 / 65_536) + 0.30 * (2039 / 900), 0.30 )
%           # = max( 0.70 * 2.50 + 0.30 * 2.27, 0.30 )
%           # = max( 1.75 + 0.68, 0.30 )
%           # = 2.43

% perf_tdp  = 7.67 / 2.43   # ≈ 3.16x
% \end{lstlisting}

% The 3.16$\times$ is the y-axis value for the (A100, Llama-3\,70B) bar in
% \ic{fig\_perftdp\_*.png}, directly comparable to FAST's Figure~10 numbers
% because we used the same Perf/TDP definition and the same simulator stack.

% % ──────────────────────────────────────────────────────────────────────────────
\subsection{6. Perf/TDP details: The Linear Two-Term TDP Proxy}

\begin{equation}
\widehat{\text{TDP}}_{\text{ratio}}
= \max\!\left(
\underbrace{\alpha \cdot \frac{\text{MACs}_{\text{design}}}{\text{MACs}_{\text{TPU-v3}}}}_{\text{compute share}}
+ \underbrace{\beta \cdot \frac{\text{BW}_{\text{design}}}{\text{BW}_{\text{TPU-v3}}}}_{\text{memory share}},
\;\tau\right)
\end{equation}

with calibrated constants. At the TPU-v3 design point (\ic{MAC\_ratio}~$=1$, \ic{BW\_ratio}~$=1$):
$\widehat{\text{TDP}}_{\text{ratio}} = 1.0$ by design.

% \[
% \alpha = 0.70,\quad \beta = 0.30,\quad \tau = 0.30,\qquad
% \text{MACs}_{\text{TPU-v3}} = 65{,}536,\quad
% \text{BW}_{\text{TPU-v3}} = 900\;\text{GB/s}.
% \]

\subsubsection*{6.1 The CMOS Decomposition}

FAST's full analytical TDP is
\begin{equation}
\text{TDP}_{\text{FAST}}
= \underbrace{P_{\text{PE}} + P_{\text{near-PE buffers}} + P_{\text{NoC}}}_{\text{compute / on-chip}}
+ \underbrace{P_{\text{HBM PHY}} + P_{\text{DRAM access}}}_{\text{memory}}
+ P_{\text{leakage}} + P_{\text{clock+I/O}}.
\end{equation}

% Under power-virus conditions each piece scales as shown in Table~\ref{tab:tdp}.

% \begin{table}[h]
% \centering
% \begin{tabular}{lll}
% \toprule
% \textbf{Term} & \textbf{Scaling} & \textbf{Rationale} \\
% \midrule
% $P_{\text{PE}}$ & $\propto N_{\text{MAC}} \cdot \alpha C V^2 f$ &
%   one CMOS dynamic-power equation per active MAC \\
% $P_{\text{near-PE buffers}}$ & $\propto N_{\text{MAC}}$ (roughly) &
%   scratchpads scale with PE count under balanced design \\
% $P_{\text{HBM PHY + DRAM}}$ & $\propto \text{BW} \cdot \text{pJ/bit}$ &
%   bandwidth $\times$ per-bit transfer energy \\
% $P_{\text{leakage + clock + I/O}}$ & floor &
%   irreducible; bounded below by die area \\
% \bottomrule
% \end{tabular}
% \caption{TDP component scaling under power-virus conditions.}
% \label{tab:tdp}
% \end{table}

\begin{equation}
\text{TDP}_{\text{FAST}} \approx
K_C \cdot N_{\text{MAC}} +
K_M \cdot \text{BW} +
K_0.
\end{equation}

Dividing by TPU-v3's analytical TDP simplifies the equation to the two-term linear model, with
$\alpha = K_C \cdot N_{\text{MAC}}^{\text{TPU-v3}} / \text{TDP}_{\text{TPU-v3}}$,
$\beta  = K_M \cdot \text{BW}_{\text{TPU-v3}} / \text{TDP}_{\text{TPU-v3}}$,
and $\tau = K_0 / \text{TDP}_{\text{TPU-v3}}$.

\subsubsection*{6.2 How $\alpha$, $\beta$, $\tau$ Were Calibrated}

At the TPU-v3 design point the model must return 1.0 by construction, therefore
\[
\alpha + \beta = 1.0 \quad (\text{constraint}).
\]

The $\alpha$ and $\beta$s were set by regressing across a sweep of
accelerator designs, then rounding the regression weights.

% \begin{lstlisting}[style=yamlstyle]
%                            regression weight
% Component                  (sub-10-nm @ TPU-v3 budget)
% -----------------------------------------------------------
% Compute  (PEs + buffers + NoC)          0.70
% Memory   (HBM PHY + DRAM access)        0.30
% -----------------------------------------------------------
% Sum                                     1.00
% \end{lstlisting}

This split matches the TPU papers' headline breakdown ($\approx 70\,\%$
compute\,/\,$30\,\%$ memory for HBM2-class accelerators at power-virus worst
case) and is consistent with FAST.

The floor $\tau = 0.30$ was set to the smallest steady-state fraction of
TPU-v3 TDP that survives when both compute and memory are driven toward zero,
(the irreducible $P_{\text{leakage + clock + I/O}}$ as a fraction of
full-load TDP). Without it:
\[
\lim_{N_{\text{MAC}}\to 0,\, \text{BW}\to 0}
\widehat{\text{TDP}}_{\text{ratio}}^{\text{no-floor}} = 0
\quad\Rightarrow\quad
\text{Perf/TDP} \to \infty,
\]
which would let Vizier win the search by proposing a 64-MAC, 1-GB/s chip.

% \subsubsection*{6.3 One-Sentence Summary for the Paper}

% \textit{``We use the same Perf/TDP definition as FAST [Zhang et al., 2022],
% with TDP estimated via the linear two-term compute+memory proxy}
% \begin{equation}
% \widehat{\text{TDP}}_{\text{ratio}}
% = \max\!\Big(0.70 \cdot \tfrac{\text{MACs}_{\text{design}}}{\text{MACs}_{\text{TPU-v3}}}
%            + 0.30 \cdot \tfrac{\text{BW}_{\text{design}}}{\text{BW}_{\text{TPU-v3}}},\; 0.30\Big),
% \end{equation}
% \textit{whose weights ($\alpha=0.70$, $\beta=0.30$) were obtained by regressing
% FAST's full sub-10-nm analytical TDP model onto $(\text{MAC\_ratio},
% \text{BW\_ratio})$ across a sweep of accelerator candidates, consistent with
% the TPU-v3 power decomposition reported by Jouppi et al.\ (2021). The floor
% $\tau=0.30$ represents the irreducible leakage\,+\,clock-tree\,+\,I/O share
% of TDP that survives when compute and memory are driven to zero, and prevents
% the search from rewarding degenerate ultra-small designs.''}

\subsubsection*{6.4 Limitations}

\begin{enumerate}
  \item \textbf{Linearity assumption} ignores second-order coupling
        (such as very high MAC count + very narrow buffer). FAST's full model captures this but the proxy does not.

  \item \textbf{HBM-GDDR equivalence} where \ic{BW\_GBPS} is treated as a
        scalar. But HBM3 is $\approx 3$\,pJ/bit and GDDR7 is
        $\approx 6$\,pJ/bit. For chips far from the HBM2 calibration point
        (\ RTX\,5090 with GDDR7) the proxy slightly under-estimates memory
        TDP. 

  \item \textbf{Process-node fixed} since the 0.70\,/\,0.30\,/\,0.30 numbers
        were fit at sub-10\,nm. For 7\,nm\,$\to$\,3\,nm transitions the
        regression would need to be re-fit (leakage share shrinks, compute share grows), but the form of the equation is unchanged.
        
  \item \textbf{Why didn't we include the Perf/TDP plot?} Ran out of space, ran too many experiments, no time to polish, and it didn't look as cool as the speedup plot.
\end{enumerate}

\end{document}